% 原书分章 PDF 第 8—10 页：第 23.5 节；重复条件编号按原印保留。
\section{定理 23.4 条件的讨论}

\noindent\textbf{条件 23.3 和 23.3。}这是试图估计式 (23.3) 所定义的 $\kappa$ 分位数 $\eta_0$ 时使用的标准条件。只要能很好地估计 $\bar Q_{b,0}$，这些条件就能保证：收集大量数据能够使我们很好地估计随机变量 $\bar Q_{b,0}$ 的 $\kappa$ 分位数。关于“很好地估计”$\bar Q_{b,0}$ 的含义，参见引理 23.1。当 $W$ 连续且 $\eta_0\ne0$ 时，预期这些条件成立似乎是合理的，因为我们假定 $\bar Q_{b,0}$ 不会退化于任意的（由 $\kappa$ 决定的）点 $\eta_0$。

\noindent\textbf{条件 23.3。}如果 $\tau_0>0$，那么 23.3 由 23.3 蕴含。如果 $\tau_0=0$，那么 23.3 类似于非例外分布假设，即 $W$ 的某个分层中不存在处理效应的概率为零。由于事先不知道 $\tau_0$，我们要求稍强一点的条件，即 $W$ 的某个分层中出现任何一个特定的小处理效应的概率为零。注意，这个条件并不排除处理效应很小，例如对所有 $\tau>0$ 都有 $P_0(|\bar Q_{b,0}(W)|<\tau)>0$；它排除的是存在一个序列 $\tau_m\downarrow0$，使得 $P_0(\bar Q_{b,0}(W)=\tau_m)>0$ 无穷多次成立。因此，这个条件似乎并没有真正比非例外分布假设更强。如果担心这类例外分布，我们建议把 Luedtke 与 van der Laan（2016b）的方法推广到 R-C 情形。

\noindent\textbf{条件 23.4。}这个条件假定，每个协变量分层中的个体都有一个合理的（至少为 $\delta>0$ 的）处理概率。

\noindent\textbf{条件 23.4。}这个条件要求，$g_0$ 的估计尊重每个协变量分层都有合理处理概率这一事实。

\noindent\textbf{条件 23.4。}如果 $R_{10}(d_n,P_n^*)=o_{P_0}(n^{-1/2})$ 且 $R_{20}(d_n)=o_{P_0}(n^{-1/2})$，那么这个条件成立。$R_{10}(d_n,P_n^*)$ 具有典型的双重稳健项的形式：只要 $g_0$ 或 $\bar Q_0$ 之一估计得好，它就很小；而且它是二阶的，即如果 $g_0$ 与 $\bar Q_0$ 都估计得好，我们就可以期望 $R_{10}(d_n,P_n^*)=o_{P_0}(n^{-1/2})$。使用 Cauchy--Schwarz 不等式，可以用这两个量的 $L^2(P_0)$ 收敛速率之积来给出该余项的上界。如果 $g_0$ 已知，可以取 $g_n=g_0$，这时该项为零。

保证 $R_{20}(d_n)=o_{P_0}(n^{-1/2})$ 需要多做一些工作，不过仍会证明这是一个合理的条件。对于某些 $C_0,\alpha>0$，我们将使用下列间隔假设（margin assumption）：
\begin{equation}
 P_0\bigl(0<|\bar Q_{b,0}-\tau_0|\leq t\bigr)\leq C_0t^\alpha
 \quad\text{对所有 }t>0.
\tag{23.5}
\end{equation}
这个间隔假设类似于 Audibert 与 Tsybakov（2007）使用的假设。以下结果把 $R_{20}(d_n)$ 的收敛速率与 $\bar Q_{b,n}-\tau_n$ 向 $\bar Q_{b,0}-\tau_0$ 的收敛速率联系起来。

\begin{bookresult}{定理}{23.5}
如果对于某些 $C_0,\alpha>0$，式 (23.5) 成立，那么对于某个常数 $C>0$，有
\[
 |R_{20}(d_n)|\leq C\min\left\{
 \begin{gathered}
 \bigl\|(\bar Q_{b,n}-\tau_n)-(\bar Q_{b,0}-\tau_0)\bigr\|_{2,P_0}^{2(1+\alpha)/(2+\alpha)},\\
 \bigl\|(\bar Q_{b,n}-\tau_n)-(\bar Q_{b,0}-\tau_0)\bigr\|_{\infty,P_0}^{1+\alpha}
 \end{gathered}\right\}.
\]
\end{bookresult}

该结果的证明类似于定理 22.5，故省略。如果 $S_0$ 在 $\tau_0$ 处具有有限导数，如 (23.3) 所给出的那样，那么可以取 $\alpha=1$。这时，上述定理表明，如果
$\bigl\|(\bar Q_{b,n}-\tau_n)-(\bar Q_{b,0}-\tau_0)\bigr\|_{2,P_0}$
为 $o_{P_0}(n^{-3/8})$，或者
$\bigl\|(\bar Q_{b,n}-\tau_n)-(\bar Q_{b,0}-\tau_0)\bigr\|_{\infty,P_0}$
为 $o_{P_0}(n^{-1/4})$，则 $R_{20}(d_n)=o_{P_0}(n^{-1/2})$。

\noindent\textbf{条件 23.4。}这是一个温和的一致性条件；$d_n$、$g_n$ 和 $\bar Q_n^*$ 分别在 $L^2(P_0)$ 中一致于 $d_0$、$g_0$ 和 $\bar Q_0$，即可推出该条件。注意，在 23.4 下，初始（尚未波动的）估计 $\bar Q_n^0$ 对 $\bar Q_0$ 的一致性，会蕴含 $\bar Q_n^*$ 对 $\bar Q_0$ 的一致性，因为在这种情形下 $\epsilon_n\to0$ 依概率成立，从而 $\|\bar Q_n^*-\bar Q_n^0\|_{2,P_0}\to0$ 依概率成立。

\noindent\textbf{条件 23.4。}这个条件限制了 $d_0$、$g_0$ 和 $\bar Q_0$ 的估计量所能具有的数据自适应程度。关于何种条件可使“$d_0$、$g_0$ 和 $\bar Q_0$ 的估计属于 Donsker 类”蕴含“$D(d_n,\tau_0,P_n^*)$ 属于 Donsker 类”，参见 van der Vaart 与 Wellner（1996）第 2.10 节。van der Laan 与 Luedtke（2015）用交叉验证 TMLE 估计价值函数，避开了这个条件；使用这种技术，在这里也可以避开这一条件。

\noindent\textbf{条件 23.4。}对于任意分布 $P$ 和函数 $f:\mathcal O\to\mathbb R$，采用记号 $Pf=\int f(o)\,dP(o)$，则有
\[
\begin{aligned}
 P_nD(d_n,\tau_0,P_n^*)&=P_nD_1(d_n,P_n^*)+P_nD_2(d_n,P_n^*)\\
 &\quad-\tau_0\left(\frac1n\sum_{i=1}^n\E_{P_U}[d_n(U,w_i)]-\kappa\right).
\end{aligned}
\]
由 TMLE 算法的波动步骤，第一项为零；由于 $P_n^*$ 使用 $W$ 的经验分布作为 $W$ 的边际分布，右侧第二项也为零。如果 $\tau_0=0$，第三项显然为零，因此假定 $\tau_0>0$。结合式 (23.4) 以及 $d_n$ 是替代估计量这一事实可知，只要 $\tau_n>0$ 以趋于 1 的概率成立，第三项就以趋于 1 的概率为零。当然，如果 $\tau_n\to\tau_0>0$ 依概率成立，就会出现这种情形；引理 23.1 为此给出了充分条件。
